\documentclass[11pt]{article}

% ----------------------------------------------------------------------
% Packages
% ----------------------------------------------------------------------
\usepackage[utf8]{inputenc}
\usepackage[T1]{fontenc}
\usepackage{amsmath,amssymb}
\usepackage{booktabs}
\usepackage[margin=1in]{geometry}
\usepackage{xcolor}
\usepackage{graphicx}
\graphicspath{{visuals/}}
\usepackage[colorlinks=true,linkcolor=blue,citecolor=blue,urlcolor=blue]{hyperref}
\usepackage{enumitem}

% ----------------------------------------------------------------------
% Title
% ----------------------------------------------------------------------
\title{Reinforcement Learning on the Taxi Environment:\\
       Tabular Q-learning and a DQN Comparison}
\author{Ana Abou-Arraj Awad \and Asier Ugarteche Perez}
\date{\today}

\begin{document}
\maketitle

\begin{abstract}
We solve the Gymnasium \texttt{Taxi} environment with \textbf{tabular Q-learning} and
study how hyperparameter choices affect performance. After a manually tuned baseline, we
compare three automated tuning strategies --- grid search, random search, and Hyperband
(via Optuna). All four configurations ultimately solve the task with a $100\%$ success
rate, which illustrates a key point: on a small, deterministic environment like Taxi,
hyperparameters influence \emph{how fast} the agent learns more than \emph{how well} it
ends up performing. We close with a section on the \textbf{Deep Q-Network (DQN)} approach,
which replaces the Q-table with a neural network --- including a multi-seed hyperparameter
study (network size, regularization, and the RL knobs) and a replay-buffer $\times$
target-network ablation --- and discuss when that added machinery is warranted. As an optional
extension we also apply the same DQN to Atari Pong from raw pixels, swapping the tabular-style MLP
for a convolutional network and vanilla DQN for Double DQN.
\end{abstract}

% ======================================================================
\section{The Taxi environment}
% ======================================================================
Taxi is a $5\times5$ grid world. The agent drives a taxi to a passenger, picks them up,
drives to a destination, and drops them off. The environment has:
\begin{itemize}[noitemsep]
  \item \textbf{500 states} --- $25$ taxi positions $\times\;5$ passenger locations
        (four pickup spots plus ``inside the taxi'') $\times\;4$ destinations;
  \item \textbf{6 actions} --- South, North, East, West, Pickup, Dropoff;
  \item \textbf{rewards} --- $-1$ per step, $+20$ for a correct dropoff, $-10$ for an
        illegal pickup/dropoff.
\end{itemize}
Crucially, Taxi is \emph{deterministic}: a trained agent should deliver the passenger every
single time. With only 500 discrete states, the problem is small enough that a lookup
table over $(\text{state}, \text{action})$ pairs is a natural and sufficient representation
--- which is exactly what tabular Q-learning uses.

% ======================================================================
\section{Tabular Q-learning (without DQN)}
% ======================================================================

\subsection{The algorithm}
Q-learning stores one value $Q(s,a)$ per state--action pair in a
$500 \times 6$ table, initialised to zero. At each step the agent takes an
$\epsilon$-greedy action and updates the table using the Bellman rule:
\begin{equation}
  Q(s,a) \;\leftarrow\; Q(s,a) \;+\; \alpha\,\Bigl[\,r + \gamma \max_{a'} Q(s',a') - Q(s,a)\,\Bigr].
  \label{eq:bellman}
\end{equation}
Exploration is annealed over training with an exponential schedule,
\begin{equation}
  \epsilon(\text{ep}) \;=\; \epsilon_{\min} + (\epsilon_{\max} - \epsilon_{\min})\,
  e^{-\,\texttt{decay\_rate}\,\cdot\,\text{ep}} ,
  \label{eq:epsilon}
\end{equation}
so the agent explores heavily early and increasingly exploits its learned policy later.

\subsection{Hyperparameters}
\label{sec:hyperparameters}
The behaviour of tabular Q-learning is governed entirely by a handful of hyperparameters.
Table~\ref{tab:hyperparams} lists each one, its symbol in Eqs.~\eqref{eq:bellman}--\eqref{eq:epsilon},
its role, and the value used in the manual baseline (Part~1).

\begin{table}[h]
\centering
\begin{tabular}{@{}llcl@{}}
\toprule
\textbf{Hyperparameter} & \textbf{Symbol} & \textbf{Manual value} & \textbf{Role} \\
\midrule
\texttt{learning\_rate} & $\alpha$            & $0.7$    & step size of the Bellman update \\
\texttt{gamma}          & $\gamma$            & $0.95$   & discount on future reward \\
\texttt{total\_episodes}& ---                 & $25{,}000$ & number of training episodes \\
\texttt{max\_steps}     & ---                 & $200$    & step cap per episode \\
\texttt{max\_epsilon}   & $\epsilon_{\max}$   & $1.0$    & initial exploration rate \\
\texttt{min\_epsilon}   & $\epsilon_{\min}$   & $0.01$   & exploration floor \\
\texttt{decay\_rate}    & ---                 & $0.0005$ & speed of $\epsilon$ decay \\
\bottomrule
\end{tabular}
\caption{Q-learning hyperparameters and their manual-baseline values.}
\label{tab:hyperparams}
\end{table}

\noindent\textbf{How each one matters.}
\begin{itemize}[noitemsep]
  \item \texttt{learning\_rate} ($\alpha$): how strongly each new experience overwrites the
        current estimate. High values learn fast but can be unstable; low values are smoother
        but slower.
  \item \texttt{gamma} ($\gamma$): how much the agent values future reward. Close to $1$ makes
        it far-sighted, which suits Taxi because the big $+20$ reward only arrives at the end
        of a successful trip.
  \item \texttt{min\_epsilon} and \texttt{decay\_rate}: together they shape the
        exploration-to-exploitation transition. A \emph{low} \texttt{min\_epsilon} lets the
        agent commit to its learned policy once it knows enough; \texttt{decay\_rate} sets how
        quickly it gets there. These two turned out to be the most influential hyperparameters
        (Section~\ref{sec:comparison}).
  \item \texttt{total\_episodes} and \texttt{max\_steps}: the training budget and the per-episode
        horizon. $200$ steps is comfortably above the $\sim$13 steps an optimal policy needs.
\end{itemize}

\noindent\textbf{Evaluation note.} Throughout, the trained Q-table is evaluated with a
\textbf{pure greedy policy} ($\epsilon = 0$), averaged over $5$ seeds $\times$ $200$ episodes
$= 1{,}000$ evaluation episodes per method. Because Taxi is deterministic, the only source of
variance is the randomised initial state; averaging over seeds and episodes gives a stable,
reproducible score.

\subsection{Part 1 --- Manual training}
With the values in Table~\ref{tab:hyperparams}, the agent trained for $25{,}000$ episodes.
The learning curve (Figure~\ref{fig:manual-curve}) climbs steeply in the first few thousand
episodes and then flattens once a solid policy is found. The evaluated baseline:
\[
  \text{avg reward } 7.93, \qquad \text{avg steps } 13.1, \qquad \text{success } 100\%.
\]

\begin{figure}[ht]
\centering
\includegraphics[width=0.85\linewidth]{manual_learning_curve.png}
\caption{Manual-baseline learning curve (first $8{,}000$ episodes, where learning happens):
training reward as a $100$-episode rolling average (left axis, blue) and the exploration rate
$\epsilon$ (right axis, orange dotted). Reward climbs steeply over the first few thousand
episodes and then flattens as $\epsilon$ decays and the agent shifts from exploring to
exploiting. Code:
\href{https://github.com/asugar13/rl-taxi-pong/blob/main/QLearning_Taxi.ipynb}{\texttt{QLearning\_Taxi.ipynb}}, Part~1.}
\label{fig:manual-curve}
\end{figure}

\subsection{Part 2 --- Grid search}
Grid search trains every combination of a fixed set of values: \texttt{learning\_rate}
$\in\{0.1, 0.3, 0.5, 0.7, 0.9\}$, \texttt{gamma} $\in\{0.85, 0.95, 0.99\}$, \texttt{decay\_rate}
$\in\{0.0005, 0.002\}$, and \texttt{min\_epsilon} $\in\{0.01, 0.05\}$ --- $60$ combinations in
total. Each was trained for $15{,}000$ episodes and scored by the mean reward over the last
$1{,}000$ episodes. It is exhaustive but constrained to the predefined grid points.
The best configuration scored $7.62$:
\[
  \texttt{learning\_rate}=0.9,\quad \texttt{gamma}=0.99,\quad
  \texttt{decay\_rate}=0.0005,\quad \texttt{min\_epsilon}=0.01 .
\]
The dominant pattern was \texttt{min\_epsilon}: every \texttt{min\_epsilon}$=0.01$ panel
outperformed its \texttt{min\_epsilon}$=0.05$ counterpart by roughly $2$--$2.5$ reward points,
regardless of the other hyperparameters. Once \texttt{min\_epsilon} is low, the environment is
tolerant of \texttt{learning\_rate}, \texttt{gamma}, and \texttt{decay\_rate} --- no single one
of them produces a sharp peak.

\subsection{Part 3 --- Random search}
Instead of a fixed grid, random search samples $30$ configurations from continuous ranges:
\texttt{learning\_rate} $\sim U(0.1, 0.9)$, \texttt{gamma} $\sim U(0.80, 0.999)$,
\texttt{decay\_rate} $\sim U(0.0001, 0.005)$, \texttt{min\_epsilon} $\sim U(0.001, 0.1)$. With
the same number of trials it covers the space better than grid because it is not pinned to
predetermined points. The best trial scored $7.92$:
\[
  \texttt{learning\_rate}=0.646,\quad \texttt{gamma}=0.828,\quad
  \texttt{decay\_rate}=0.00108,\quad \texttt{min\_epsilon}=0.00173 .
\]
The scatter-matrix analysis again singled out \texttt{min\_epsilon} as by far the strongest
predictor: every top-$10$ trial had \texttt{min\_epsilon} below $0.03$, while every bottom-$10$
trial had it above $0.07$ --- the two groups barely overlap. \texttt{decay\_rate} shows a much
weaker secondary trend; \texttt{learning\_rate} and \texttt{gamma} show no clean individual
pattern.

\subsection{Part 4 --- Hyperband with Optuna}
Grid and random search waste compute by training every configuration for the full budget,
regardless of how it is doing. \textbf{Hyperband} runs the search like a tournament: it starts
many configurations on a small budget, eliminates the worst, and gives more episodes only to
the survivors. We implement this with \textbf{Optuna} (a TPE sampler combined with its
\texttt{HyperbandPruner}) over the same continuous ranges as random search, for $40$ trials ---
of which $31$ were pruned early by Hyperband, with only $9$ running to completion.
The best trial scored $7.74$:
\[
  \texttt{learning\_rate}=0.888,\quad \texttt{gamma}=0.871,\quad
  \texttt{decay\_rate}=0.0049,\quad \texttt{min\_epsilon}=0.0021 .
\]

Optuna's own hyperparameter-importance estimate (fANOVA over the completed trials)
confirms the ordering seen in the grid and random searches (Figure~\ref{fig:optuna-importance}):
\texttt{min\_epsilon} dominates ($\approx 0.52$), with the $\epsilon$ \texttt{decay\_rate} a clear
second ($\approx 0.31$) and \texttt{learning\_rate} and \texttt{gamma} contributing little
($\approx 0.09$ each). In other words, how exploration is shut down --- both its floor and how
fast it gets there --- matters far more than the learning rate or discount on this deterministic
environment.

\begin{figure}[ht]
\centering
\includegraphics[width=0.85\linewidth]{optuna_importance.png}
\caption{Optuna hyperparameter importance (fANOVA) over the $40$-trial Hyperband study.
\texttt{min\_epsilon} is by far the strongest predictor of final score, followed by the
$\epsilon$ \texttt{decay\_rate}; \texttt{learning\_rate} and \texttt{gamma} are near-negligible.
Code:
\href{https://github.com/asugar13/rl-taxi-pong/blob/main/QLearning_Taxi.ipynb}{\texttt{QLearning\_Taxi.ipynb}}, Part~4.}
\label{fig:optuna-importance}
\end{figure}

\subsection{Part 5 --- Final comparison}
\label{sec:comparison}
The best configuration from each method was retrained for the full $25{,}000$ episodes and
evaluated identically. Table~\ref{tab:search} summarises the best configurations and their
search scores; Table~\ref{tab:final} reports the final evaluation.

Plotting the four retrained runs side by side (Figure~\ref{fig:all-curves}) shows that the methods
differ in \emph{how fast} they learn far more than in where they end up: all four climb to the same
plateau, but the Optuna/Hyperband config converges fastest --- it reaches near-optimal reward within
roughly the first thousand episodes, whereas the manual and grid configs take several thousand. The
reason is visible in Table~\ref{tab:search}: Optuna's config pairs a high learning rate with a much
faster $\epsilon$ decay ($0.0049$ vs.\ $0.0005$), so it stops exploring and commits to its policy
early. Manual and grid share the slow decay $0.0005$ and their curves almost coincide. This is the
same lesson as before --- on a deterministic environment with a unique optimal policy, tuning buys
faster convergence rather than a better final policy.

\begin{figure}[ht]
\centering
\includegraphics[width=0.85\linewidth]{all_learning_curves.png}
\caption{Learning curves (reward, $200$-episode rolling average) for the best configuration of each
method, retrained for $25{,}000$ episodes and zoomed to the first $8{,}000$ where the differences
appear. All four reach the same plateau; Optuna/Hyperband (green) gets there fastest thanks to its
faster $\epsilon$ decay, while manual (blue) and grid (orange) share the slow decay and nearly
overlap. Code:
\href{https://github.com/asugar13/rl-taxi-pong/blob/main/QLearning_Taxi.ipynb}{\texttt{QLearning\_Taxi.ipynb}}, Part~5.}
\label{fig:all-curves}
\end{figure}

\begin{table}[h]
\centering
\begin{tabular}{@{}lcccccc@{}}
\toprule
\textbf{Method} & $\alpha$ & $\gamma$ & \texttt{decay\_rate} & \texttt{min\_epsilon} & \textbf{Search score} \\
\midrule
Manual          & $0.70$ & $0.95$ & $0.0005$  & $0.0100$  & --- \\
Grid search     & $0.90$ & $0.99$ & $0.0005$  & $0.0100$  & $7.62$ \\
Random search   & $0.646$& $0.828$& $0.00108$ & $0.00173$ & $7.92$ \\
Optuna/Hyperband& $0.888$& $0.871$& $0.0049$  & $0.0021$  & $7.74$ \\
\bottomrule
\end{tabular}
\caption{Best configuration found by each method (search score = mean reward over the last
$1{,}000$ of $15{,}000$ training episodes for grid/random/Optuna; manual uses $25{,}000$ episodes).}
\label{tab:search}
\end{table}

\begin{table}[h]
\centering
\begin{tabular}{@{}lccc@{}}
\toprule
\textbf{Method} & \textbf{Avg reward} & \textbf{Avg steps} & \textbf{Success rate} \\
\midrule
Manual           & $\mathbf{7.93}$ & $\mathbf{13.1}$ & $100\%$ \\
Grid search      & $\mathbf{7.93}$ & $\mathbf{13.1}$ & $100\%$ \\
Random search    & $7.92$ & $13.1$ & $100\%$ \\
Optuna/Hyperband & $7.83$ & $13.2$ & $100\%$ \\
\bottomrule
\end{tabular}
\caption{Final evaluation after retraining each best config for $25{,}000$ episodes
(pure greedy, $\epsilon = 0$, averaged over $5$ seeds $\times$ $200$ episodes).}
\label{tab:final}
\end{table}

\noindent\textbf{Analysis.} Three observations stand out:
\begin{enumerate}[noitemsep]
  \item \textbf{All four methods solve Taxi} ($100\%$ success, ${\sim}13.1$--$13.2$ steps), and
        their final rewards fall within $0.10$ of each other. On a deterministic environment,
        hyperparameters affect \emph{how fast} the agent converges far more than the final
        policy quality, since the optimal policy is unique and easily representable in a table.
  \item \textbf{During the search phase}, random search found the highest-scoring configuration
        ($7.92$), followed by grid ($7.62$) and Optuna ($7.74$). After full retraining for
        $25{,}000$ episodes, those differences washed out: all four methods are statistically
        indistinguishable ($7.83$--$7.93$).
  \item \textbf{Compute cost} is where the methods differ most meaningfully. Grid search
        consumed $900{,}000$ episodes ($60$ combinations $\times$ $15{,}000$). Random search used
        half that --- $450{,}000$ episodes --- for a result within $0.01$ reward of grid search's
        best. Optuna/Hyperband used only ${\sim}164{,}000$ episodes, pruning $31$ of $40$ trials
        early, and still landed within $0.10$ of the top performer: roughly a \textbf{quarter} of
        grid search's compute for a comparable result.
\end{enumerate}
The consistent thread across all searches is that \texttt{min\_epsilon} is by far the most
influential hyperparameter: a low exploration floor lets the agent commit fully to its learned
policy. \texttt{decay\_rate} has a weaker secondary effect; \texttt{learning\_rate} and
\texttt{gamma} show no clear independent pattern once \texttt{min\_epsilon} is controlled for.

% ======================================================================
\section{Deep Q-Network (DQN) approach}
\label{sec:dqn}
% ======================================================================
The natural question is what changes if we replace the Q-table with a neural network --- a
\textbf{Deep Q-Network}. The algorithm is identical (the Bellman update of
Eq.~\eqref{eq:bellman}); only the \emph{representation} of $Q$ changes.

\subsection{From a Q-table to a function}
A Q-table stores one number per $(s,a)$ pair --- a lookup. A DQN instead \emph{computes}
$Q(s,a)$ with a neural network. Because the network cannot consume a raw state ID (the integer
$499$ is not ``larger'' than $12$ in any meaningful sense), each Taxi state is converted to a
\textbf{one-hot vector} of length $500$ before being fed to the network. This makes the DQN on
Taxi close to a tabular method in disguise: with one-hot inputs the network learns essentially
one independent column of weights per state, so there is little generalisation across states.

\smallskip\noindent\textbf{Two ways to apply the same Bellman target.} The methods also
\emph{update} differently. Tabular Q-learning edits one cell in place using the
temporal-difference (TD) error --- the bracketed quantity in Eq.~\eqref{eq:bellman},
$\delta = r + \gamma \max_{a'} Q(s',a') - Q(s,a)$, scaled by the learning rate $\alpha$. A neural
network cannot edit individual outputs, so a DQN instead forms a fixed \textbf{bootstrap target}
from the (frozen) target network,
\[
  y \;=\; r + \gamma \max_{a'} Q_{\theta^-}(s',a'),
\]
and \textbf{minimises a regression loss} between its prediction and that target by gradient descent,
\[
  \mathcal{L}(\theta) \;=\; \bigl(\,Q_{\theta}(s,a) - y\,\bigr)^2
  \qquad\text{(Huber / smooth-}L_1\text{ in practice)},
\]
where $\theta$ is the online (policy) network and $\theta^-$ the target network. So the explicit
subtraction of Eq.~\eqref{eq:bellman} becomes a squared (or Huber) loss the optimiser descends ---
the same Bellman principle, different machinery (in-place table edit vs.\ backprop on a function
approximator).

\subsection{Architecture and key components}
The DQN solution\footnote{Full code and outputs: \href{https://github.com/asugar13/rl-taxi-pong/blob/main/DQN_Taxi.ipynb}{\texttt{DQN\_Taxi.ipynb}}.} uses:
\begin{itemize}[noitemsep]
  \item a small multilayer perceptron, $500 \rightarrow 64 \rightarrow 64 \rightarrow 6$
        (one-hot input, one Q-value per action);
  \item an \textbf{experience replay buffer} (capacity $50{,}000$) that stores past transitions
        and samples random minibatches of $64$, breaking the correlation between consecutive
        steps. Sampling reads transitions without removing them, so the same experience may be
        reused in later updates; only when the buffer reaches capacity does adding a new
        transition discard the oldest one;
  \item a \textbf{target network} --- a slowly-updated copy of the policy network, synced every
        $250$ steps --- used to compute a stable Bellman target;
  \item the Huber (smooth-$L_1$) loss with gradient clipping, optimised by Adam.
\end{itemize}
These two stabilisers (replay buffer and target network) are the machinery that distinguishes
DQN from plain online Q-learning with a function approximator.

\subsection{Replay warm-up and per-step learning}
The replay threshold is measured in \emph{transitions}, not episodes. Although training runs for
$500$ episodes, the replay warm-up ends as soon as $500$ individual environment steps have been
stored, normally within the first few episodes because each episode can last up to $200$ steps.
Before that threshold, the agent continues to act and populate the buffer, but no backpropagation
or Q-value update occurs. Action selection is still called ``$\epsilon$-greedy'': with probability
$\epsilon$ the action is random, while with probability $1-\epsilon$ it is the argmax of the
policy network. During warm-up that network is still randomly initialised, however, so its rare
``exploit'' actions are not yet informed by learning.

After each environment step the new transition $(s,a,r,s',d)$ is added first. Once the buffer
contains at least $500$ transitions, \texttt{train\_frequency = 1} means that every subsequent
step is followed by one optimizer update from a randomly sampled batch of $64$ stored
transitions. These samples do not empty the replay buffer. The target network is updated
separately by copying the policy-network weights every $250$ environment steps.

\subsection{Baseline configuration and result}
\label{sec:dqn-baseline}
The baseline full DQN uses $\gamma = 0.99$, learning rate $10^{-3}$, batch size $64$,
replay capacity $50{,}000$, replay warm-up $500$ transitions, one optimizer update per
environment step after warm-up, and a target-update period of $250$ steps. Its exponential
exploration schedule approaches the floor $0.05$ rather than reaching it within the run:
$\epsilon$ is approximately $0.999$ at episode~1, $0.868$ at episode~100, and $0.499$ at
episode~500. This deliberately slow decay keeps collecting varied experience; it does not mean
that the final policy is evaluated with $50\%$ random actions. Training lasts $500$ episodes,
whereas evaluation is fully greedy ($\epsilon = 0$). The resulting agent reaches:
\[
  \text{avg reward } 8.04, \qquad \text{avg steps } 13.0, \qquad \text{success } 100\%.
\]

\subsection{Hyperparameter study}
\label{sec:dqn-hpo}
Mirroring the tabular search, we tune the DQN in two stages (coordinate descent), each trained
over \textbf{multiple seeds} ($\{7,8,9\}$) and reported as \textbf{mean $\pm$ std}. The
multi-seed protocol is deliberate: with a single seed the \emph{ranking} of near-equal variants
can flip from run to run, so averaging over seeds (and reporting the spread) is what makes a
``best'' claim defensible.

\paragraph{Stage~1 --- network size and regularization.}
Holding the full DQN fixed, we vary the architecture --- a single $32$-unit hidden layer, the
$64$\nobreakdash-$64$ baseline, and a wider $128$\nobreakdash-$128$ --- and the regularization
(L2 weight decay $10^{-4}$, dropout $0.1$), each trained for $400$ episodes over three seeds
(Table~\ref{tab:dqn-net}; learning curves in Figure~\ref{fig:dqn-net}).

\begin{table}[h]
\centering
\begin{tabular}{@{}lrrr@{}}
\toprule
\textbf{Variant} & \textbf{Params} & \textbf{Success \% (mean)} & \textbf{Reward (mean $\pm$ std)} \\
\midrule
$128$-$128$ + weight decay & $81{,}414$ & $\mathbf{100.0}$ & $\mathbf{7.99 \pm 0.04}$ \\
wide $128$-$128$           & $81{,}414$ & $94.7$ & $-3.11 \pm 15.71$ \\
baseline $64$-$64$         & $36{,}614$ & $85.0$ & $-23.15 \pm 23.36$ \\
$64$-$64$ + weight decay   & $36{,}614$ & $81.7$ & $-29.93 \pm 30.49$ \\
shallow $32$               & $16{,}230$ & $63.7$ & $-67.31 \pm 28.05$ \\
$64$-$64$ + dropout $0.1$  & $36{,}614$ & $0.0$  & $-200.0 \pm 0.0$ \\
$128$-$128$ + dropout $0.1$& $81{,}414$ & $0.0$  & $-200.0 \pm 0.0$ \\
\bottomrule
\end{tabular}
\caption{Network / regularization variants --- 3 seeds, 400 training episodes, greedy eval.}
\label{tab:dqn-net}
\end{table}

\noindent The standout is \textbf{$128$-$128$ with weight decay}: it is the \emph{only} variant to
fully solve Taxi --- $\mathbf{100\%}$ success in the optimal ${\sim}13$ steps, with essentially zero
variance ($\pm 0.04$ over seeds). Weight decay's effect is \textbf{architecture-dependent}: it
\emph{rescues} the high-capacity wide net (from $94.7\%$ to a perfect score) yet \emph{hurts} the
small $64$-$64$ net ($85.0\% \to 81.7\%$) --- consistent with the larger network being the one
prone to the over-fitting/instability that regularization tames. Network size alone helps (wide
$128$-$128$ beats the $64$-$64$ baseline beats the shallow $32$), but only \emph{with} weight decay
does it become rock-solid. \textbf{Dropout is catastrophic} on both sizes ($0\%$) --- injecting
noise into the Q-estimates breaks DQN. Finally, the winning $128$-$128$+weight-decay matches the
single-seed Part-1 figure (Section~\ref{sec:dqn-baseline}: $8.04$ reward, $100\%$) but now
\emph{across three seeds with a $\pm0.04$ spread}: unlike the noisier variants (whose multi-seed
averages sit well below that figure --- a reminder that one lucky seed can flatter), this
configuration is genuinely robust.

\begin{figure}[ht]
\centering
\includegraphics[width=0.85\linewidth]{fig_net_search.png}
\caption{Stage~1 learning curves --- training reward (mean over 3 seeds, $50$-episode moving
average) for the seven network/regularization variants. \mbox{$128$-$128$ + weight decay} converges
highest (to the optimal return); the two dropout variants flatline at the failure floor. Code:
\href{https://github.com/asugar13/rl-taxi-pong/blob/main/DQN_Taxi.ipynb}{\texttt{DQN\_Taxi.ipynb}}, Part~2.1.}
\label{fig:dqn-net}
\end{figure}

\paragraph{Stage~2 --- RL hyperparameters.}
Table~\ref{tab:dqn-rl-space} lists the knobs we search --- what each controls, the range it was
sampled from, and its default (Part~1) value.

\begin{table}[h]
\centering
\small
\begin{tabular}{@{}l p{5.4cm} l l@{}}
\toprule
\textbf{Knob} & \textbf{What it does} & \textbf{Search range} & \textbf{Default} \\
\midrule
\texttt{lr}               & Adam step size --- how fast the weights move    & $3{\times}10^{-4}$--$5{\times}10^{-3}$ & $10^{-3}$ \\
\texttt{gamma} ($\gamma$) & discount on future reward (far-sightedness)     & $0.95$--$0.999$                       & $0.99$ \\
\texttt{decay}            & rate at which exploration $\epsilon$ decays     & $5{\times}10^{-4}$--$4{\times}10^{-3}$ & $1.5{\times}10^{-3}$ \\
\texttt{batch\_size}      & transitions sampled per gradient step           & $\{32,\,64,\,128\}$                   & $64$ \\
\texttt{train\_freq}      & one gradient step every $N$ environment steps   & $\{1,\,2,\,4\}$                       & $1$ \\
\texttt{memory}           & replay-buffer capacity (transitions kept)       & $2{,}000$--$50{,}000$ (log)           & $50{,}000$ \\
\texttt{update}           & target-network update mechanism                 & hard / soft                           & hard \\
\texttt{target\_update}   & \emph{(hard)} env steps between target syncs    & $100$--$1000$                         & $250$ \\
\texttt{tau} ($\tau$)     & \emph{(soft)} fraction the target shifts toward the policy each step & $0.001$--$0.02$ (log) & --- \\
\bottomrule
\end{tabular}
\caption{Stage~2 search space --- the RL hyperparameters, what each controls, the range each was
sampled from, and its default (Part~1) value.}
\label{tab:dqn-rl-space}
\end{table}

A random search of $16$ configurations --- over the learning knobs (learning rate, $\gamma$,
$\epsilon$-decay), the replay/optimisation settings (batch size, learning frequency), and the
\textbf{target-update mechanism} (a hard periodic copy vs.\ a soft Polyak update
$\theta^- \leftarrow \tau\,\theta + (1-\tau)\,\theta^-$) --- found that \emph{none} beats the
defaults: no random trial outscored the default recipe in the single-seed scan, and the
hard-vs-soft choice in particular made no difference, the default hard update already being best. Combined with
Stage~1, the recommended configuration is therefore $128$\nobreakdash-$128$ \textbf{with weight
decay} and the \emph{default} RL hyperparameters, scoring $\mathbf{7.99 \pm 0.04}$ reward at
$\mathbf{100\%}$ success (the standout row of Table~\ref{tab:dqn-net}) --- Taxi solved in the
optimal ${\sim}13$ steps. The RL knobs are thus effectively \emph{neutral}: the standard recipe is
already best, and what lifts the agent to a perfect score is the \emph{weight decay} from
Stage~1, not the learning knobs.

\noindent The search \emph{failures} are themselves informative. The lowest-scoring configurations
share a $1{:}4$ learning frequency (a gradient step only every fourth environment step --- hence
undertrained at the $250$-episode scan) and/or a tiny replay buffer ($\sim$2--3k transitions, prone
to forgetting): what breaks the agent is \emph{undertraining and forgetting}, not the
learning-rate/$\gamma$/$\epsilon$-decay knobs, which show no clean trend. The hard-vs-soft target
choice in particular shows no pattern --- both appear at the top \emph{and} the bottom of the
ranking. The defaults already avoid both failure modes (train every step, $50$k buffer), which is
why they are robust and nothing beats them. (Tendencies only: the scan is single-seed and varies
every knob at once.) The full 16-trial search table is in the
\href{https://github.com/asugar13/rl-taxi-pong/blob/main/DQN_Taxi.ipynb}{notebook} (Part~2.2).

\subsection{Ablation: replay buffer and target network}
\label{sec:dqn-ablation}
To isolate what each algorithm component contributes, we train all four corners of the
\emph{replay buffer} $\times$ \emph{target network} $2\times 2$ --- both, buffer-only,
target-only, and neither (online, batch size $1$) --- again over multiple seeds, reported as
mean $\pm$ std. Importantly, the network and hyperparameters are held at a \textbf{neutral
baseline} rather than the tuned best from Section~\ref{sec:dqn-hpo}: tuned settings become
\emph{co-adapted} to the full system, so reusing them would unfairly handicap the ablated corners
(e.g.\ a learning rate that is stable only \emph{because} of the target network would blow up the
target-only corner). A shared neutral reference keeps the comparison fair.

\paragraph{Results.}
The outcome is unambiguous (Table~\ref{tab:dqn-ablation}, Figure~\ref{fig:dqn-ablation}). The \textbf{replay buffer is decisive}:
without it (\emph{target-only} and \emph{neither}) the agent never learns the sparse reward and
fails completely ($0\%$ success, $-200$ reward). With it (\emph{both} and \emph{buffer-only}) it
solves Taxi at $87$--$91\%$. The \textbf{target network adds only a modest gain} --- \emph{both}
($91.3\%$) over \emph{buffer-only} ($87.0\%$) --- so it is a useful refinement, not a necessity.
On near-tabular Taxi, the replay buffer carries most of the load.

\begin{table}[h]
\centering
\begin{tabular}{@{}lccrr@{}}
\toprule
\textbf{Variant} & \textbf{Replay} & \textbf{Target} & \textbf{Success \% (mean $\pm$ std)} & \textbf{Reward (mean $\pm$ std)} \\
\midrule
both        & yes & yes & $91.3 \pm 12.3$ & $-10.06 \pm 25.54$ \\
buffer only & yes & no  & $87.0 \pm 16.3$ & $-18.91 \pm 33.77$ \\
target only & no  & yes & $0.0 \pm 0.0$   & $-200.0 \pm 0.0$ \\
neither     & no  & no  & $0.0 \pm 0.0$   & $-200.0 \pm 0.0$ \\
\bottomrule
\end{tabular}
\caption{Replay-buffer $\times$ target-network ablation --- 3 seeds, 500 episodes, greedy eval.}
\label{tab:dqn-ablation}
\end{table}

\begin{figure}[ht]
\centering
\includegraphics[width=0.85\linewidth]{fig_ablation.png}
\caption{Ablation learning curves --- training reward (mean over 3 seeds, $100$-episode moving
average) for the four replay\,$\times$\,target corners. The two buffer-equipped variants
(\emph{both}, \emph{buffer only}) climb out of the failure floor; \emph{target only} and
\emph{neither} flatline. Code:
\href{https://github.com/asugar13/rl-taxi-pong/blob/main/DQN_Taxi.ipynb}{\texttt{DQN\_Taxi.ipynb}}, Part~3.}
\label{fig:dqn-ablation}
\end{figure}

\subsection{Tabular vs.\ DQN}
\begin{table}[h]
\centering
\begin{tabular}{@{}lll@{}}
\toprule
 & \textbf{Tabular Q-learning} & \textbf{DQN} \\
\midrule
$Q$ representation & $500 \times 6$ table        & neural net (one-hot $\rightarrow$ MLP) \\
State input        & integer ID (table index)    & one-hot vector of length $500$ \\
Stabilisers        & none needed                 & replay buffer $+$ target network \\
Training episodes  & $25{,}000$                  & $500$ \\
Eval reward        & $\sim 7.93$ (greedy)        & $8.04$ (greedy) \\
Eval steps         & $\sim 13.1$                 & $13.0$ \\
Success rate       & $100\%$                     & $100\%$ \\
\bottomrule
\end{tabular}
\caption{Tabular Q-learning versus DQN on Taxi.}
\label{tab:tabvsdqn}
\end{table}

\noindent Both approaches are evaluated with a pure greedy policy, so the scores are directly
comparable. The small remaining gap ($8.04$ vs.\ $7.93$) is within run-to-run noise; both
fully solve Taxi in $\sim$$13$ steps with a $100\%$ success rate. The real takeaway is that,
for a 500-state problem, a Q-table is the simpler and more natural tool; DQN is introduced
here as a \emph{learning exercise} for how one would scale to state spaces too large to
tabulate, where a function approximator becomes mandatory.

% ======================================================================
\section{Conclusion}
% ======================================================================
Tabular Q-learning solves Taxi reliably, and on this small, deterministic problem the choice of
hyperparameters affects convergence speed more than final quality --- every tuning method we
tried reached a $100\%$ success rate. The most influential knobs were \texttt{min\_epsilon} and
\texttt{decay\_rate}, which control how decisively the agent shifts from exploration to
exploitation. The DQN reaches the same behaviour through a neural network plus a replay buffer
and a target network; on Taxi that is more machinery than the problem demands, but it is the
exact recipe one needs when the state space grows beyond what a table can hold.
A multi-seed hyperparameter study and a replay-buffer $\times$ target-network ablation sharpen the
picture. The standout configuration is a \emph{larger} network \emph{with weight decay} --- the only one to
fully solve Taxi ($100\%$, ${\sim}13$ steps, near-zero variance); weight decay rescued the wide net
while it hurt the small one, and dropout broke both. The RL hyperparameters --- including batch
size, learning frequency, and a hard-vs-soft target update --- proved close to neutral, no sampled
config beating the defaults. The ablation, however, is decisive: the \emph{replay buffer} is the
load-bearing component (without it the agent never learns the sparse reward), while the target
network is only a modest refinement.

% ======================================================================
\section{Optional extension: DQN on Pong from pixels}
\label{sec:pong}
% ======================================================================
As an optional extension we apply the \emph{same} DQN to Atari \textbf{Pong}, learned directly from
raw screen pixels.\footnote{Code: \href{https://github.com/asugar13/rl-taxi-pong/blob/main/DQN_Pong.ipynb}{\texttt{DQN\_Pong.ipynb}}.}
The learning algorithm barely changes --- replay buffer, target network, $\epsilon$-greedy, and a
bootstrapped Bellman target all carry over, the lone refinement being a \emph{Double}-DQN target
(detailed below) --- while the \emph{representation}, the \emph{scale}, and the \emph{engineering} change
substantially.
Table~\ref{tab:taxi-vs-pong} summarises the jump from Taxi.

\begin{table}[h]
\centering
\small
\begin{tabular}{@{}l p{4.4cm} p{5.8cm}@{}}
\toprule
 & \textbf{Taxi DQN} & \textbf{Pong DQN} \\
\midrule
Observation         & state index (one-hot of $500$)      & $4$ stacked $84\times84$ greyscale frames \\
One obs.\ Markov?   & yes --- the index is sufficient     & no --- stack $4$ frames (see below) \\
Q-network           & MLP ($128$--$128$)                  & CNN (3 conv $\to$ $512$ dense), ``Nature DQN'' \\
Input size          & $500$ discrete states               & ${\sim}28$k pixels per state $(4{\times}84{\times}84)$ \\
Action space        & $6$ (all used)                      & $6\to3$ (NOOP / up / down) \\
Reward              & native ($-10$ to $+20$)             & clipped to $\{-1,0,+1\}$ \\
DQN target          & vanilla ($\max$)                    & \textbf{Double} DQN \\
Replay buffer       & $50$k transitions                   & $100$k \texttt{uint8} images (${\sim}6$\,GB) \\
$\epsilon$ schedule & decayed over episodes               & annealed over $200$k frames \\
Target sync         & every $250$ steps                   & every $2$k frames \\
Optimiser           & Adam, \texttt{lr}$=10^{-3}$         & Adam, \texttt{lr}$=10^{-4}$ \\
Scale / compute     & ${\sim}$hundreds of episodes, CPU   & ${\sim}1$--$2$M frames, ${\sim}1.5$--$2$\,h on GPU/MPS \\
\bottomrule
\end{tabular}
\caption{From Taxi to Pong: the learning \emph{algorithm} is identical, but the representation
(pixels vs.\ an index), the network (CNN vs.\ MLP), and the scale all change. The $\epsilon$,
action-space, and target-sync rows turned out to decide whether Pong learned at all
(Section~\ref{sec:pong-debug}).}
\label{tab:taxi-vs-pong}
\end{table}

\subsection{Observation, state, and the four-frame Markov window}
In Taxi the environment hands us a state \emph{index}: a sufficient statistic that already satisfies
the \textbf{Markov property} --- the next state and reward depend only on it and the chosen action. In
Pong the observation is an \textbf{image}, and a single frame is \emph{not} Markovian: from one still
you can read the ball's position but not its \emph{direction} or \emph{speed}, so the optimal action is
genuinely undefined. We restore the Markov property the standard way --- by \textbf{stacking the last
four frames} as the state,
\[
  s_t \;=\; \big(x_{t-3},\,x_{t-2},\,x_{t-1},\,x_t\big)\in\{0,\dots,255\}^{4\times84\times84},
\]
so that velocity (and even acceleration) becomes recoverable from differences across the stack. That
four-frame window is precisely the piece Taxi got for free from its index; in Pong we have to construct
it. The pipeline that produces each frame $x_t$: convert to greyscale, resize to $84\times84$, skip
(and max-pool) $4$ emulator frames per action, then stack $4$ such frames. A \textbf{convolutional}
Q-network (three conv layers, then a $512$-unit dense layer, then one Q-value per action) reads the
stack; rewards are clipped to $\{-1,0,+1\}$, the $100$k \texttt{uint8} buffer holds memory at
${\sim}6$\,GB, and checkpointing lets a disconnected GPU session resume.

\subsection{Vanilla vs.\ Double DQN}
The two notebooks also differ in \emph{one line} of the target computation. Both keep an online
network $\theta$ and a target network $\theta^-$; the question is who \emph{selects} and who
\emph{evaluates} the next action. \textbf{Vanilla} DQN (Taxi) lets the target network do both,
fused into a single max:
\[
  y \;=\; r + \gamma \max_{a'} Q_{\theta^-}(s', a') .
\]
\textbf{Double} DQN (Pong) splits the two roles --- the online network \emph{selects} the action,
the target network \emph{evaluates} it:
\[
  a^\star = \arg\max_{a'} Q_{\theta}(s', a'),
  \qquad
  y \;=\; r + \gamma\, Q_{\theta^-}(s', a^\star) .
\]
\noindent\textbf{Why.} The $\max$ over noisy Q-estimates is biased \emph{upward}: it preferentially
selects whichever action got lucky positive noise, so the max of the estimates sits above the max
of the true values, and because DQN bootstraps, that overestimation feeds back on itself. Splitting
selection from evaluation \emph{decorrelates} the noise --- if the online net over-rates an action
by chance, the frozen target net usually does not over-rate that same action, so the inflation
largely cancels (van Hasselt, Guez \& Silver, 2016). The cost is essentially nil: both variants
already maintain the two networks, so going Double is a one-line change to the target. Taxi uses
vanilla because it is near-tabular and overestimation is harmless there; Pong's noisier pixel
function-approximation makes the bias bite harder, so Double is the natural choice.

\subsection{Results and final configuration}
\label{sec:pong-debug}
Our first full run never learned: at two million frames the score sat at the random-play baseline
(${\approx}-21$) from start to finish, and greedy play lost every point. Instrumenting the agent showed
this was not a bug but a \emph{configuration} that could not learn --- with $\epsilon$ annealed over the
first $1$M of $2$M frames and the full $6$-action set, the policy collapsed onto a single action
(near-identical Q-values everywhere) and never recovered. Two changes fix it: anneal $\epsilon$
\emph{fast} (over $200$k frames) so the agent starts exploiting early, and expose only the three actions
Pong needs (NOOP / up / down); a more frequent target sync ($2$k frames) helps too.

With the corrected configuration (Table~\ref{tab:pong-config}) the \emph{same} network and update code
learn Pong from raw pixels. The average score climbs off $-21$ within ${\sim}200$--$300$k frames, reaches
break-even near $1$M, and a winning score by ${\sim}2$M\footnote{Recording of the trained agent in play:
\href{https://github.com/asugar13/rl-taxi-pong/blob/main/pong.mp4}{\texttt{pong.mp4}}.}. In the saved run,
the final training average over the last $20$ games is $7.45$; greedy evaluation over $5$ held-out episodes
scores $14.6$ on average, and the recorded game scores $16.0$. As an earlier checkpoint, the $3$-action
agent already averages ${\approx}-10$ at just $600$k frames --- a competitive paddle that returns most
balls --- against the original recipe's flat $-21$ across the entire run.

The learning curve for this run is shown in Figure~\ref{fig:pong-score}. The per-episode score (faint)
is very noisy --- individual Pong games swing widely --- but the $20$-episode moving average tells the
story cleanly: it sits on the $-21$ floor for the first ${\sim}200$ episodes while the buffer fills and
$\epsilon$ is still high, then climbs steadily, crosses break-even around episode ${\sim}490$, and
finishes averaging ${\approx}+7$. In other words the agent goes from losing every point to consistently
beating the built-in opponent --- learned end-to-end from raw pixels.

\begin{figure}[ht]
\centering
\includegraphics[width=0.72\linewidth]{pong_score.png}
\caption{Pong-from-pixels learning curve for the corrected configuration (Table~\ref{tab:pong-config}):
per-episode score (faint blue) and its $20$-episode moving average (orange). The average rises off the
$-21$ random-play floor, crosses break-even ($0$, dashed), and ends positive --- the agent learns to win
from raw pixels. Code:
\href{https://github.com/asugar13/rl-taxi-pong/blob/main/DQN_Pong.ipynb}{\texttt{DQN\_Pong.ipynb}}, Step~10.}
\label{fig:pong-score}
\end{figure}

\begin{table}[h]
\centering
\small
\begin{tabular}{@{}ll@{}}
\toprule
\textbf{Setting} & \textbf{Value} \\
\midrule
Observation         & $4\times84\times84$ greyscale stack (frame-skip $4$) \\
Action set          & $3$ --- NOOP / up / down \\
Reward              & clipped to $\{-1,0,+1\}$ \\
Q-network           & ``Nature'' CNN: 3 conv $\to$ $512$-unit dense \\
Target              & Double DQN, synced every $2$k frames \\
\midrule
Optimiser           & Adam, \texttt{lr}$=10^{-4}$ \\
Discount $\gamma$   & $0.99$ \\
Loss / grad clip    & Huber, gradient-norm $\le 10$ \\
Batch size          & $32$ \\
Replay buffer       & $100$k \texttt{uint8} transitions (${\sim}6$\,GB) \\
Learning starts     & $10$k frames \\
Train frequency     & $1$ gradient step per $4$ frames \\
$\epsilon$ schedule & $1.0 \to 0.02$, linear over $200$k frames \\
Total frames        & $2{\times}10^{6}$ \\
\bottomrule
\end{tabular}
\caption{Final Pong-from-pixels configuration. The settings that turned the original flat-lining run
into a learning one are the $3$-action set, the fast $\epsilon$-decay, and the $2$k-frame target sync;
every other line is the same DQN recipe used for Taxi, scaled up.}
\label{tab:pong-config}
\end{table}

In short, the same DQN framework that solved Taxi also solves Pong from pixels: the core machinery ---
experience replay, a target network, $\epsilon$-greedy exploration, and the bootstrapped Bellman target
--- carries over. The changes are the representation (a CNN over a four-frame stack), one
\emph{algorithmic} refinement --- the \emph{Double}-DQN form of the target (online net selects, target
net evaluates) in place of Taxi's vanilla $\max$ --- and the handful of training settings that decide
whether exploration is ever converted into a winning policy.

\end{document}
